\section{Experimental Methodology}

The experiments are organized around the fixed-budget placement claim. The
first layer asks whether moving a host-level background budget toward tenants
currently carrying write pressure improves their write path without extra
capacity. The second layer narrows the mechanism: runtime actuation tests
deployable control, stronger baselines test fixed skew and local autotuning,
score instances test the adaptation signal, ranking replay and score-family
perturbations test robustness, cgroups test a lower-layer control point, and
the rate-limiter coverage microstudy checks the maintenance-I/O knob. The third
layer reports workload coverage, tier-wise collateral cost, liveness, and trace
qualification boundaries. Budget conservation is enforced within each
experiment family. The epoch-level and public-trace qualification runs use an
11/7/3\,MB/s, 112\,MB/s contract; the continuous runtime family uses an
11/6/1\,MB/s, 96\,MB/s contract. These numbers differ because the harnesses
expose different actuation surfaces, but every SAKI comparison holds the
relevant family budget fixed against its baselines.

\subsection{Workloads}

The epoch-level anchor is \epochdrift. It uses 16 co-located tenants, three
demand tiers, eight epochs, and gradual movement of a four-tenant
high-pressure set. Tenants are prefilled before the measured epochs.

The continuous anchor is \runtimedrift. It also uses 16 tenants and three
tiers, but runs tenants as long-lived embedded RocksDB processes for eight
20-second windows. The high-pressure set moves by one tenant per window, so the
experiment tests online adaptation at a rate the controller can plausibly
follow. The continuous anchor uses 11/6/1\,MB/s high/mid/low budgets with a
4/8/4 split, preserving a 96\,MB/s aggregate against an all-mid static
baseline. A window-length sensitivity bracket repeats the same anchor at 10,
20, and 40 seconds for three static/online seed pairs per point, holding the
160-second duration and all per-tenant budgets fixed. A score-family
robustness ablation perturbs the demand-mode instance along anchor weight and a
uniform residual scale ($\times 0.5$, $\times 2$, and $\times 0$ meaning
anchor-only), with three seeds per row paired with paper-static a/b/c; results
appear in Section~\ref{sec:evaluation}. We also replay all five continuous main
online traces offline to measure the ranking margin and the top-\(H\)
invariance of nearby score instances without launching additional RocksDB
runs. A separate 360-second longer-run is a sanity check for obvious
instability.

We also run a local sensitivity sweep around \epochdrift, changing drift
speed (\texttt{epochs\_per\_stage}=3 and 1, versus 2) and high write pressure
(70k and 130k writes/s, versus 100k). Each new point starts with two
\texttt{static}/SAKI screening trials; stress-side variants are promoted to
five trials if a pre-committed rule holds (high P99 improves by at least 2\%,
high throughput by at least 5\%, overlap $\ge$ 2.5/4, static does not hit
capacity-boundary).

Around \runtimedrift we run a controlled workload matrix that perturbs one
axis at a time: \valuefourk uses 4\,KB values, \readthreex triples per-tier
point-read QPS, \fastdrift moves the high-pressure set at 1.5 tenants per
window, and \tenantseight uses a smaller 8-tenant cluster with two
high-pressure tenants. Per-tenant budgets are preserved so the same-budget
contract holds; for \tenantseight the aggregate scales linearly. The matrix
uses the paper SAKI policy without changing the score instance. Every variant is
promoted from a two-trial screen to five trials unless either policy has a
failed tenant, and the overlap threshold is rescaled to each variant's
high-pressure tenant count.
The 20-second window, tier budgets, and 4/8/4 split are mechanism-isolation
settings, not portable constants. The 4/8/4 split is chosen because the
high-budget tier's surplus exactly matches the low-budget tier's savings,
preserving the same aggregate budget within each experiment family;
Section~\ref{sec:saki-design} states the general conservation rule.
\fastdrift and \tenantseight are local sensitivity, not a parameter search.

\subsection{Policies}

\texttt{static} is the main fair baseline: every tenant receives the same share
of the aggregate background budget for that experiment family. SAKI is the
online policy: it assigns high, mid, and low tiers from previous-interval
observations while preserving that aggregate budget. The main continuous SAKI
policy uses the demand-aware reference score instance.

\staticbiased is a frozen-skew baseline for the epoch workload. It uses the
same high/mid/low budget tiers, but freezes the initial-stage assignment for
all later epochs. It models an operator who pre-tunes a reasonable non-equal
allocation and then does not adapt as the high-pressure set drifts.

\staticautotuned is the strict same-budget built-in \rocksdb baseline. It uses
exactly the same per-tenant ceilings as \texttt{static}. In the epoch
configuration this is two background jobs and 7\,MB/s per tenant, for a 32-job
and 112\,MB/s aggregate budget across 16 tenants. Each tenant enables
\texttt{--rate\_limiter\_auto\_tuned=1}, so adaptation stays local to that
RocksDB instance and remains within the configured per-instance range. This
baseline asks whether per-instance local rate adaptation can replace the outer
host-level placement decision; we keep the ceilings identical to \texttt{static}.
For the sensitivity sweep, \staticautotuned is added only to the promoted
higher-pressure point.

\texttt{oracle\_tiered} is a diagnostic reference using true tier information.
\texttt{unlimited} is a capacity reference, not a same-budget competitor.

\subsection{Platform}
\label{sec:platform}

All experiments run on a single Linux machine with 80 logical cores (Ubuntu,
kernel 6.17). Storage is a local SSD. Each RocksDB tenant writes to its own
data directory on the same device. The epoch-level harness uses \texttt{db\_bench}
8.9.1 from unmodified Ubuntu \texttt{.deb} archives; the continuous harness
links against the same RocksDB version. No other major I/O-intensive processes
run during measured trials.

\subsection{Metrics and Aggregation}

Primary metrics are write P99/P999 latency and write throughput for
high-pressure tenants, the foreground symptoms SAKI is meant to improve by
moving background service. Total throughput shows whether the host collapses
aggregate work under the fixed budget. Tier-wise latency and throughput expose
the redistribution cost that a targeted policy can impose on lower-pressure
tenants. The continuous rate limiter is background-I/O control: it shapes
RocksDB flush/compaction I/O, so foreground changes reflect maintenance
placement rather than direct sleeps on application reads or writes.

Epoch-level efficiency reports compaction write bytes. In continuous runs SAKI
changes foreground completion, so the main compaction efficiency metric is
normalized bytes/write; absolute compaction bytes are reported but not the main
efficiency claim. Epoch runs execute comparable foreground work per trial, so
absolute background bytes are directly interpretable; continuous runs may
complete more foreground writes under SAKI and need per-write normalization.

Adaptation is measured by high-set overlap (high-pressure tenants assigned the
high-budget tier after warmup) and adaptation lag in control windows. Ranking
robustness is measured by replaying the previous-window observations that fed
each online decision and comparing the resulting top-\(H\) set across score
instances using exact-match count and Jaccard similarity. Liveness is measured
by failed tenant processes. The main epoch and continuous results aggregate
five independent trials. Percent changes are computed relative to static within
each trial before aggregation. The paper reports mean, range, standard
deviation, and confidence intervals. Sensitivity rows with two trials report a
mean and range; five-trial rows report 95\% CI half-widths. Raw and summarized
artifacts are recorded under \texttt{remote-results/}, and figures use real
result files from the figure data manifest.

\subsection{Scope}

The methodology supports a constrained claim: fixed-budget online reallocation
in two controlled RocksDB drift regimes. The built-in local rate-limiter
baseline is covered under the same aggregate budget because it tests whether
tenant-local adaptation can replace the host-level rank-and-assign loop.
OS/cgroup I/O controls are included for a different reason: they can enforce
weights or throttles below RocksDB, but they do not observe LSM debt signals or
decide which engine should receive compaction service. We include a 5-trial
aligned cgroup-v2 baseline at the same scale as the epoch main result (16
tenants, 8 epochs, 240k keys/tenant, 100k writes/high tier) as supplemental
mechanism evidence, not as a repeated main baseline. The controller inside the
cgroup harness uses a simplified pressure score for cross-mechanism comparison;
it is not the main SAKI controller used in the primary RocksDB experiments.
